% ============================================================
% §1  Introduction
% ============================================================
\section{Introduction}
\label{sec:introduction}

Bayesian optimization (BO) is the standard approach to optimizing
expensive black-box functions: objectives that are costly to evaluate,
return no gradients, and may be non-convex or multimodal
\citep{frazier2018tutorial,shahriari2016taking}.
Applications range from hyperparameter tuning in machine learning to
aerodynamic shape optimization, where each function evaluation can take
minutes to hours of computation or physical experiment.
BO sequentially builds a probabilistic surrogate model—typically a
Gaussian process (GP)—and uses an \emph{acquisition function} to decide
where to evaluate next.
The acquisition function must balance \emph{exploitation} (querying
where the surrogate predicts high values) against \emph{exploration}
(querying where the surrogate is uncertain, to reduce the risk of
missing a better region).

A natural question is whether the GP's uncertainty estimate alone is
sufficient to drive sample efficiency.
Purely uncertainty-driven sampling—selecting the point of maximum
posterior standard deviation—uses the GP's calibration to target
information gain, without any regard for whether the point is predicted
to have a high objective value.
If uncertainty-driven sampling outperforms blind random search, it would
suggest that the GP's uncertainty model, rather than the exploitation
signal, is the primary engine of BO's advantage.
If, on the other hand, it performs no better than random, it would
implicate exploitation as the necessary component.
Classical acquisition functions such as Expected Improvement (EI) and
Upper Confidence Bound (UCB) combine both signals; understanding which
component matters more is fundamental to designing and interpreting BO
algorithms.

I address this question through a systematic benchmark.
I evaluate four strategies—EI, UCB, Uncertainty (maximum posterior
standard deviation), and Random—on six benchmark problems spanning
dimensions 2 to 10 and two problem families: classical synthetic
functions and physics-based engineering test functions.
All strategies share the same GP surrogate and BO loop, so any
performance difference is attributable solely to the acquisition
function.
I use a Friedman–Nemenyi non-parametric statistical framework
\citep{demsar2006statistical} to draw conclusions robust to
non-Gaussian regret distributions.

\paragraph{Contributions.}
\begin{enumerate}[nosep]
  \item \textbf{Systematic 4-strategy benchmark.}
    I run 4 strategies $\times$ 6 problems $\times$ 10 seeds $= 240$
    BO runs under identical conditions and analyze results with a
    non-parametric test (Friedman + Nemenyi) rather than mean-only
    comparisons.
    The analysis reveals two statistically separated equivalence classes:
    balanced strategies \{EI, UCB\} and pure-exploration strategies
    \{Uncertainty, Random\}.

  \item \textbf{GP normalization is necessary on multi-scale engineering
    problems.}
    An experimental comparison on Piston (7D, inputs spanning
    approximately seven orders of magnitude) shows that omitting GP input
    normalization degrades BO to random-search performance, while the
    Random baseline is unaffected.
    This isolates the fault to the surrogate and validates the choice
    of explicit, bounds-based normalization.

  \item \textbf{Exploitation, not uncertainty, is the critical ingredient.}
    The statistical equivalence of Uncertainty and Random (Nemenyi
    $p = 0.974$) provides a controlled decomposition: the GP's
    posterior variance, used without the posterior mean, provides no
    measurable advantage over blind Sobol sampling.
    Exploitation of the GP mean is the necessary condition for
    sample-efficient BO.
\end{enumerate}

The remainder of the paper is organized as follows.
Section~\ref{sec:background} reviews the BO framework, the GP surrogate,
and related acquisition function families.
Section~\ref{sec:methods} details the experimental design, surrogate
implementation, and statistical methodology.
Section~\ref{sec:results} presents all benchmark results.
Section~\ref{sec:discussion} interprets the findings and discusses
limitations.
Section~\ref{sec:conclusion} concludes.

% ============================================================
% §2  Background
% ============================================================
\section{Background}
\label{sec:background}

\subsection{The Bayesian Optimization Loop}
\label{sec:bo-background}

Bayesian optimization solves $\max_{x \in \mathcal{X}} f(x)$ when $f$
is a black-box function that is expensive to evaluate, returns no
gradients, and may be multimodal \citep{frazier2018tutorial}.
The algorithm proceeds in two alternating steps.
First, a \emph{surrogate model} is fitted to all observations
$\{(x_i, y_i)\}_{i=1}^n$; it provides a posterior distribution over
$f$ with a calibrated mean $\mu(x)$ and variance $\sigma^2(x)$.
Second, an \emph{acquisition function} $\alpha(x)$ converts the
posterior into a scalar score at each candidate point; the next
evaluation location is $x_{n+1} = \argmax_{x \in \mathcal{X}} \alpha(x)$.
Solving the inner maximization is cheap because $\alpha$ is an
analytical function of the posterior.

BO is best suited to problems where: (i) each evaluation of $f$ is
expensive (minutes to hours), (ii) the dimensionality is low-to-moderate,
and (iii) gradients of $f$ are unavailable
\citep{frazier2018tutorial,shahriari2016taking}.

\subsection{Gaussian Process Surrogate}
\label{sec:gp-background}

A Gaussian process (GP) is a distribution over functions fully specified
by a mean function $m(x)$ and a positive-definite kernel
$k(x, x')$ \citep{frazier2018tutorial}.
After observing $n$ points, the GP posterior at a new location $x_*$
is Gaussian with
\[
  \mu(x_*) = k_*^\top (K + \sigma_n^2 I)^{-1} \mathbf{y},
  \qquad
  \sigma^2(x_*) = k(x_*, x_*) - k_*^\top (K + \sigma_n^2 I)^{-1} k_*,
\]
where $K_{ij} = k(x_i, x_j)$, $k_* = [k(x_i, x_*)]_i$, and
$\sigma_n^2$ is the noise variance.
The key computational cost is the $n\times n$ matrix inversion (or
Cholesky decomposition), which is $\mathcal{O}(n^3)$ and limits exact
GP inference to $n \lesssim 10^3$ observations \citep{shahriari2016taking}.

The kernel encodes smoothness assumptions about $f$.
I use the radial basis function (RBF) kernel with
\emph{automatic relevance determination} (ARD): one independent
lengthscale $\ell_j$ per input dimension $j$.
ARD allows the GP to learn which dimensions carry most of the variation
in $f$, effectively performing soft variable selection.
Kernel hyperparameters are fitted by maximizing the exact marginal
log-likelihood at each BO iteration.

\subsection{Acquisition Functions}
\label{sec:acq-background}

Acquisition functions trade off exploitation of high-mean regions
against exploration of high-variance regions.
All classical acquisition functions use both $\mu(x)$ and $\sigma(x)$;
the key difference is how they weight the two signals.

\paragraph{Expected Improvement (EI).}
$\alpha_{\mathrm{EI}}(x) = \mathbb{E}[\max(f(x) - f^+, 0)]$ where
$f^+ = \max_i y_i$ is the current best observation.
Under a GP this has the closed form
$\alpha_{\mathrm{EI}}(x) = \sigma(x)[\gamma\Phi(\gamma) + \phi(\gamma)]$
with $\gamma = (\mu(x) - f^+)/\sigma(x)$ \citep{frazier2018tutorial}.
EI is positive even when $\mu(x) < f^+$, provided $\sigma(x)$ is
large—this is why EI naturally explores uncertain regions.
EI has no finite-time regret guarantee.

\paragraph{Upper Confidence Bound (UCB).}
$\alpha_{\mathrm{UCB}}(x) = \mu(x) + \sqrt{\beta}\,\sigma(x)$.
The parameter $\beta$ directly controls the exploration–exploitation
trade-off.
Under a theory-motivated schedule $\beta_t \propto \log t$, GP-UCB
achieves sub-linear cumulative regret \citep{frazier2018tutorial}.
In practice a fixed $\beta$ is common; I use $\beta = 2$.

\paragraph{Pure-exploration strategies.}
Maximizing $\sigma(x)$ alone—without any $\mu(x)$ term—corresponds to
selecting the point about which the GP is most uncertain.
This pure-exploration strategy is useful as a controlled ablation: it
uses the GP's calibration machinery but discards all exploitation signal.
Blind random search serves as the weakest baseline.
De~Ath et al.~\citeyear{deAth2019greed} give a theoretical framing:
they show EI and UCB both lie on the exploration–exploitation Pareto
front in the ($\mu$, $\sigma$) plane, whereas maximum-$\sigma$ sampling
is off this front—it is Pareto-dominated in the exploitation dimension.
Their empirical results suggest this off-front position is
consequential in higher dimensions.

\subsection{Related Work}
\label{sec:related}

\paragraph{Comparative benchmarks.}
Several large-scale BO benchmarks have compared acquisition strategies.
Shahriari et al.~\citeyear{shahriari2016taking} review EI, UCB,
Thompson Sampling, and Entropy Search across diverse tasks, and conclude
that the choice of surrogate model typically matters more than the
fine-grained choice of acquisition function—a claim my normalization
experiment instantiates concretely.
De~Ath et al.~\citeyear{deAth2019greed} compare standard acquisition
functions against $\varepsilon$-greedy methods on ten synthetic benchmarks plus
real-world tasks, finding that simple greedy strategies are competitive
with EI and UCB in high dimensions.

\paragraph{Statistical methodology.}
\Citet{demsar2006statistical} argues that comparing classifiers—or, by
extension, optimization algorithms—over multiple datasets requires
non-parametric rank-based tests rather than mean-performance
comparisons, because regret distributions are typically non-Gaussian and
the null hypothesis of equal distributions is not testable with a single
paired $t$-test across datasets.
I adopt the Friedman–Nemenyi framework he recommends for exactly this
reason.
